%mac
%\loai{Tính giá trị trung bình}



%\loai{Xác định sai số dụng cụ}


%\loai{Xác định sai số tuyệt đối}

%\loai{Xác định sai số tỉ đối khi biết kết quả đo}


%\loai{Xác định sai số tỉ đối khi biết số liệu đo}



%\loai{Viết kết quả đo chính xác}





%\loai{Xác định sai số tuyệt đối gián tiếp}


%\loai{Xác định sai số tỉ đối gián tiếp}

\begin{ex}%book
Một vật chuyển động thẳng đều với quãng đường đi được $s=15,4\pm 0,1 \ (m)$ trong thời gian $t=4,0\pm 0,1\ (s)$. Biết tốc độ của vật chuyển động thẳng đều được tính bằng công thức: $v=\dfrac{s}{t}$.
\begin{itemize}
\item[a.] Phép đo tốc độ của vật có sai số tỉ đối và sai số tuyệt đối bằng bao nhiêu? 
\item[b.] Viết kết quả đo tốc độ của vật.
\end{itemize}
\haidong{6}
\loigiai{
\begin{itemize}
\item[a.] 
\begin{itemize}
\item Sai số tỉ đối: $\delta v=\delta s+\delta t=\dfrac{\Delta s}{\overline{s}}.100\%+\dfrac{\Delta t}{\overline{t}}.100\%=3,1\%$
\item Sai số tuyệt đối: $\Delta v=\delta v.\overline{v}=\delta v.\dfrac{\overline{s}}{\overline{t}}=0,12\ m/s$ 
\end{itemize}
\item[b.] Tốc độ của vật: $\overline{v}=\dfrac{\overline{s}}{\overline{t}}=3,85\ m/s$\\
Kết quả đo tốc độ của vật: $v=3,85\pm 0,12\ m/s$. 
\end{itemize}
}
\haidong{15}
\end{ex}


